\chapter*{Conclusion Générale}
\addcontentsline{toc}{chapter}{Conclusion Générale}

Le projet \textbf{eduGenius}, réalisé au sein de l'École Polytechnique Internationale (EPI), a été une opportunité exceptionnelle de relever les défis de l'éducation moderne par la technologie. L'objectif initial était de concevoir une plateforme unifiée capable de briser les barrières de l'apprentissage à distance tout en automatisant les tâches pédagogiques répétitives.

Au terme de ce travail, nous avons abouti à une solution fonctionnelle et esthétique qui intègre :
\begin{itemize}
    \item Une gestion académique complète adaptée au système LMD.
    \item Un moteur d'intelligence artificielle performant capable de transformer des supports de cours passifs en outils d'apprentissage actifs (résumés, quiz, flashcards).
    \item Un environnement collaboratif synchrone (vidéo et chat) garantissant l'engagement des étudiants.
\end{itemize}

D'un point de vue technique, ce projet nous a permis de maîtriser l'écosystème \textbf{MERN} et d'explorer les frontières de l'IA générative. L'utilisation de \textbf{Google Gemini} et de \textbf{Daily.co} a démontré qu'il est désormais possible de bâtir des outils complexes et scalables en s'appuyant sur des services cloud de haute qualité.

Cependant, toute réalisation logicielle est perfectible. Bien que eduGenius remplisse ses objectifs primaires, plusieurs perspectives d'évolution sont envisageables :
\begin{enumerate}
    \item \textbf{Analyse prédictive} : Utiliser le \textit{Machine Learning} pour détecter précocement les risques de décrochage scolaire en analysant les scores des quiz et le taux de présence.
    \item \textbf{Gamification avancée} : Intégrer un système de badges et de classements pour stimuler la compétition saine entre étudiants.
    \item \textbf{Multimodalité} : Étendre les capacités de l'IA à l'analyse de vidéos de cours pour générer des sous-titres et des chapitres automatiques.
\end{enumerate}

En conclusion, \textbf{eduGenius} n'est pas seulement un outil de gestion, c'est un compagnon intelligent qui redonne du temps aux enseignants et de la motivation aux étudiants. Nous espérons que cette plateforme servira de base à de futures innovations pédagogiques au sein de l'EPI.
